\subsection{Pemetaan Kebutuhan terhadap Solusi}
\label{subsec:pemetaan-kebutuhan-solusi}

Hasil analisis solusi dirangkum pada Tabel \ref{tab:pemetaan-kebutuhan-solusi}. Pemetaan ini menunjukkan hubungan antara kebutuhan sistem yang ditemukan pada analisis masalah dengan komponen solusi yang akan dirancang pada bagian berikutnya.

\begin{table}[!ht]
	\centering
	\caption{Pemetaan Kebutuhan Sistem terhadap Analisis Solusi}
	\label{tab:pemetaan-kebutuhan-solusi}
	\small
	\begin{tabular}{|p{0.28\textwidth}|p{0.3\textwidth}|p{0.32\textwidth}|}
		\hline
		\textbf{Kebutuhan Sistem}                 & \textbf{Komponen Solusi}                      & \textbf{Implikasi terhadap Rancangan}                                       \\ \hline
		Kontrol akses granular                    & Macaroons dengan \textit{caveats}             & Token perlu memuat batasan aktor, kategori data, hak akses, dan waktu.      \\ \hline
		Segmentasi data \textit{off-chain}        & Pemisahan data IPFS berdasarkan kategori data & Metadata segmen data perlu dapat dicocokkan dengan \textit{caveats} token.  \\ \hline
		Delegasi berantai                         & Atenuasi \textit{offline} pada Macaroons      & Pemegang token dapat menurunkan hak akses dengan cakupan yang lebih sempit. \\ \hline
		Penegakan akses terhadap data terenkripsi & Validasi token pada \textit{server} PRE       & PRE perlu memvalidasi token sebelum menjalankan \textit{re-encryption}.     \\ \hline
	\end{tabular}
\end{table}

Dengan hasil analisis tersebut, solusi yang dirancang pada bagian berikutnya difokuskan pada empat aspek: arsitektur integrasi Macaroons ke DecMed, struktur token dan \textit{caveats}, segmentasi data \textit{off-chain}, serta alur akses dan delegasi.
